\documentclass[11pt]{article}
\usepackage[margin=1in]{geometry}
\usepackage{graphicx}
\usepackage{booktabs}
\usepackage{amsmath,amssymb}
\usepackage{enumitem}
\usepackage{hyperref}
\usepackage[table]{xcolor}
\usepackage{caption}
\usepackage{subcaption}

\definecolor{grokgreen}{HTML}{2E8B57}
\definecolor{nogrokred}{HTML}{CC3333}

\hypersetup{colorlinks=true, linkcolor=blue, citecolor=blue, urlcolor=blue}

\title{Experiment 3: Heterogeneity at the Grokking Phase Boundary}
\author{Federated Learning Grokking Project}
\date{March 2026}

\begin{document}
\maketitle

% =============================================================================
%  Experiment 3 — Heterogeneity at the Grokking Phase Boundary
% =============================================================================
\section{Experiment 3: Heterogeneity at the Phase Boundary}
\label{sec:exp3}

\subsection{Overview}

Experiment~3 investigates whether data heterogeneity shifts the grokking phase
boundary and, if so, through what mechanism.  Building on the centralized phase
boundary established in Experiment~1 ($\alpha_{\mathrm{crit}} \approx 0.2$) and
the federated baseline from Experiment~2 (FedAvg--IID matches centralized
step-for-step), we now vary the \emph{type} and \emph{severity} of
heterogeneity across clients.

The experiment has two sub-experiments:
\begin{itemize}
  \item \textbf{Exp~3a} sweeps a continuous heterogeneity knob (Dirichlet
        concentration $\alpha_{\mathrm{dir}} \in \{0.01, 0.1, 0.5, 1.0, 10.0,
        1000.0\}$) across five training fractions
        $\alpha \in \{0.2, 0.25, 0.3, 0.35, 0.5\}$ at $K{=}10$ clients, with a
        $K{=}20$ validation subset.
  \item \textbf{Exp~3b} compares structured partitions---operand-based and
        target-based---against the IID baseline at the same $\alpha$ values.
\end{itemize}
All runs use $N{=}256$, SGD with $\eta{=}50$, $\lambda{=}0$, $E{=}5$ local
epochs, full participation ($f{=}1$), and $S_{\max}{=}50{,}000$ gradient steps.
Each configuration is repeated over three seeds (42, 123, 456).


% ---------------------------------------------------------------------------
\subsection{Exp~3a: Dirichlet Heterogeneity Sweep}
\label{sec:exp3a}

\subsubsection{Phase diagram}

Figure~\ref{fig:exp3a-phase} presents the grokking phase diagram as a function
of training fraction $\alpha$ and Dirichlet concentration
$\alpha_{\mathrm{dir}}$.  The key observations are:

\begin{itemize}
  \item \textbf{The phase boundary is unchanged.}  At $\alpha{=}0.2$, grokking
        fails universally---across all six heterogeneity levels, no seed
        achieves $T_{\mathrm{grok}} < \infty$.  At $\alpha \geq 0.25$, all
        30 configurations (5~$\alpha_{\mathrm{dir}}$ values $\times$ 6
        seeds) grok to 100\% test accuracy.  Heterogeneity does not shift
        $\alpha_{\mathrm{crit}}$.

  \item \textbf{Heterogeneity delays grokking, but only near the boundary.}
        At $\alpha{=}0.25$, extreme non-IID ($\alpha_{\mathrm{dir}}{=}0.01$)
        increases $T_{\mathrm{grok}}$ from $26{,}725 \pm 2{,}048$ (IID) to
        $36{,}638 \pm 1{,}727$---a 37\% slowdown.  At $\alpha{=}0.5$, the
        effect vanishes: $T_{\mathrm{grok}}$ varies by less than 1.3\% across
        all $\alpha_{\mathrm{dir}}$ values ($7{,}638$--$7{,}738$ steps).

  \item \textbf{The slowdown saturates quickly.}  Most of the delay is
        concentrated between $\alpha_{\mathrm{dir}}{=}1000$ and
        $\alpha_{\mathrm{dir}}{=}0.1$; further increasing heterogeneity to
        $\alpha_{\mathrm{dir}}{=}0.01$ adds only modest additional delay
        (Figure~\ref{fig:exp3a-tgrok-vs-dir}).
\end{itemize}

Table~\ref{tab:exp3a} summarizes the full $\alpha \times
\alpha_{\mathrm{dir}}$ grid.  The monotonic increase in $T_{\mathrm{grok}}$
from right (IID) to left (non-IID) is visible at every $\alpha$ that groks,
but the absolute magnitude of the effect shrinks rapidly as $\alpha$ increases.

\begin{table}[htbp]
\centering
\small
\caption{\textbf{Exp~3a: Full Dirichlet sweep} ($K{=}10$, mean
$T_{\mathrm{grok}} \pm$ std over 3 seeds).  Slowdown column shows ratio
vs.\ IID ($\alpha_{\mathrm{dir}}{=}1000$).}
\label{tab:exp3a}
\begin{tabular}{lccccccr}
\toprule
$\alpha$ & $\alpha_{\mathrm{dir}}{=}0.01$ & $0.1$ & $0.5$ & $1.0$ &
$10.0$ & $1000$ (IID) & Max slow. \\
\midrule
0.20 & $\infty$ & $\infty$ & $\infty$ & $\infty$ & $\infty$ & $\infty$ & --- \\
0.25 & $36638 \pm 1727$ & $29497 \pm 3366$ & $27282 \pm 2503$ &
       $26910 \pm 2189$ & $26767 \pm 1942$ & $26725 \pm 2048$ & 1.37$\times$ \\
0.30 & $14448 \pm 831$ & $13690 \pm 595$ & $13222 \pm 493$ &
       $13190 \pm 513$ & $13087 \pm 437$ & $13072 \pm 449$ & 1.11$\times$ \\
0.35 & $10608 \pm 377$ & $10167 \pm 238$ & $9993 \pm 205$ &
       $9963 \pm 193$ & $9942 \pm 203$ & $9933 \pm 194$ & 1.07$\times$ \\
0.50 & $7738 \pm 105$ & $7680 \pm 126$ & $7638 \pm 118$ &
       $7642 \pm 120$ & $7645 \pm 113$ & $7642 \pm 120$ & 1.01$\times$ \\
\bottomrule
\end{tabular}
\end{table}

\begin{figure}[htbp]
  \centering
  \includegraphics[width=\textwidth]{figures/exp3a_phase_diagram.png}
  \caption{\textbf{Exp~3a phase diagram.}  Left: mean $T_{\mathrm{grok}}$
  (steps) as a function of $\alpha$ and $\alpha_{\mathrm{dir}}$ ($K{=}10$).
  Right: grokking success rate (fraction of seeds that grok).  Grokking
  succeeds at 100\% for all $\alpha \geq 0.25$, regardless of heterogeneity.
  $\alpha{=}0.2$ universally fails (red $\infty$ markers).}
  \label{fig:exp3a-phase}
\end{figure}

\begin{figure}[htbp]
  \centering
  \includegraphics[width=0.7\textwidth]{figures/exp3a_t_grok_vs_dir_alpha.png}
  \caption{\textbf{Grokking time vs.\ heterogeneity.}
  $T_{\mathrm{grok}}$ (mean $\pm$ std over 3 seeds) as a function of
  $\alpha_{\mathrm{dir}}$ for each training fraction.  The effect of
  heterogeneity is largest near the phase boundary ($\alpha{=}0.25$) and
  negligible in the interior ($\alpha{=}0.5$).  Note log scale on $x$-axis.}
  \label{fig:exp3a-tgrok-vs-dir}
\end{figure}


\subsubsection{$K$-validation}

Figure~\ref{fig:exp3a-kval} compares $K{=}10$ and $K{=}20$ at three
representative $(\alpha, \alpha_{\mathrm{dir}})$ settings.  Doubling the number
of clients amplifies the heterogeneity effect modestly: at $\alpha{=}0.25$ with
$\alpha_{\mathrm{dir}}{=}0.1$, $T_{\mathrm{grok}}$ increases from $29{,}497$
($K{=}10$) to $34{,}977$ ($K{=}20$), an 18.6\% further slowdown.  At
$\alpha{=}0.5$, the $K{=}10$ and $K{=}20$ curves are indistinguishable.  The
qualitative pattern is preserved: more clients slightly amplifies the delay near
the boundary but does not shift $\alpha_{\mathrm{crit}}$.

\begin{figure}[htbp]
  \centering
  \includegraphics[width=0.7\textwidth]{figures/exp3a_k_validation.png}
  \caption{\textbf{$K$-validation.}  $K{=}10$ (solid) vs.\ $K{=}20$ (dashed).
  Increasing $K$ modestly amplifies heterogeneity-induced delay near the
  boundary but does not shift the phase boundary itself.}
  \label{fig:exp3a-kval}
\end{figure}


% ---------------------------------------------------------------------------
\subsection{Exp~3b: Structured Partitions}
\label{sec:exp3b}

\subsubsection{Operand vs.\ target vs.\ IID}

A central hypothesis motivating this experiment was that operand-based
partitioning---which fragments the input space along the Fourier structure that
the model must learn---would be more disruptive than Dirichlet quantity skew.
The results \textbf{do not support this hypothesis}
(Figure~\ref{fig:exp3b-struct}, Table~\ref{tab:exp3b}).

\begin{itemize}
  \item At $\alpha{=}0.5$, all three partitions yield statistically identical
        $T_{\mathrm{grok}}$: IID~$= 7{,}642$, operand~$= 7{,}597$,
        target~$= 7{,}612$.

  \item At $\alpha{=}0.25$ (near the boundary), operand partitioning is
        actually \emph{slightly faster} than IID ($0.965\times$), while
        target partitioning is slightly slower ($1.114\times$).

  \item Both structured partitions fail at $\alpha{=}0.2$, matching the IID
        baseline---the phase boundary is unchanged.
\end{itemize}

The trajectory-level comparison (Figure~\ref{fig:exp3b-traj}) confirms that
the learning dynamics (test accuracy and IPR curves) are nearly overlapping
for all three partition types at every $\alpha$.

\begin{figure}[htbp]
  \centering
  \includegraphics[width=0.85\textwidth]{figures/exp3b_structured_partitions.png}
  \caption{\textbf{Exp~3b: Structured partition comparison.}  $T_{\mathrm{grok}}$
  by partition type for each $\alpha$ ($K{=}10$).  Operand and target
  partitions produce grokking times within noise of IID.}
  \label{fig:exp3b-struct}
\end{figure}

\begin{figure}[htbp]
  \centering
  \includegraphics[width=\textwidth]{figures/exp3b_partition_trajectories.png}
  \caption{\textbf{Structured partition trajectories} (seed~42).
  Top: test accuracy; bottom: IPR.  The three partition types produce
  nearly indistinguishable dynamics at all training fractions.}
  \label{fig:exp3b-traj}
\end{figure}

\begin{table}[htbp]
\centering
\small
\caption{\textbf{Exp~3b: Grokking time by partition type} ($K{=}10$, mean $\pm$
std over 3 seeds).  Slowdown ratio is relative to IID
($\alpha_{\mathrm{dir}}{=}1000$) from Exp~3a.}
\label{tab:exp3b}
\begin{tabular}{lccccc}
\toprule
$\alpha$ & IID ($T_{\mathrm{grok}}$) & Operand ($T_{\mathrm{grok}}$) &
Target ($T_{\mathrm{grok}}$) & Operand ratio & Target ratio \\
\midrule
0.20 & $\infty$ & $\infty$ & $\infty$ & --- & --- \\
0.25 & $26{,}725 \pm 2{,}048$ & $25{,}785 \pm 1{,}890$ &
       $29{,}783 \pm 2{,}029$ & 0.965$\times$ & 1.114$\times$ \\
0.30 & $13{,}072 \pm 449$ & $12{,}852 \pm 437$ &
       $13{,}697 \pm 324$ & 0.983$\times$ & 1.048$\times$ \\
0.35 & $9{,}933 \pm 194$ & $9{,}820 \pm 193$ &
       $10{,}093 \pm 220$ & 0.989$\times$ & 1.016$\times$ \\
0.50 & $7{,}642 \pm 120$ & $7{,}597 \pm 116$ &
       $7{,}612 \pm 95$ & 0.994$\times$ & 0.996$\times$ \\
\bottomrule
\end{tabular}
\end{table}


% ---------------------------------------------------------------------------
\subsection{Mechanistic Analysis}
\label{sec:exp3-mech}

The summary statistics establish \emph{what} happens (heterogeneity delays
grokking near the boundary but does not prevent it).  The trajectory-level
diagnostic plots reveal \emph{why}.

\subsubsection{Grokking anatomy}

Figure~\ref{fig:anatomy} overlays four metrics on a shared time axis for four
representative conditions, with vertical markers at $T_{\mathrm{onset}}$
(test accuracy first exceeds 5\%, green dotted) and $T_{\mathrm{grok}}$
(test accuracy permanently exceeds 95\%, red dashed).  The key observations:

\begin{itemize}
  \item \textbf{Heterogeneity stretches the transition, not just the onset.}
        At $\alpha{=}0.25$, the gap between $T_{\mathrm{onset}}$ and
        $T_{\mathrm{grok}}$ grows from 12{,}085 steps (IID) to 20{,}645 steps
        ($\alpha_{\mathrm{dir}}{=}0.01$)---a 71\% widening.  At $\alpha{=}0.5$,
        the gap is ${\sim}1{,}800$--$1{,}990$ steps regardless of heterogeneity.

  \item \textbf{The generalization gap peaks before $T_{\mathrm{grok}}$ and
        collapses to zero after it} (row~b).  Under non-IID, the peak is
        taller and wider, indicating that the model spends longer in the
        ``memorized but not yet generalized'' regime.

  \item \textbf{Client drift peaks just before $T_{\mathrm{grok}}$, then
        decays sharply} (row~c).  This pattern is universal across all
        conditions.  The peak drift under extreme non-IID is ${\sim}2\times$
        higher than under IID, but the post-grokking decay profile is
        identical---once the Fourier solution locks in, clients converge to
        agreement.

  \item \textbf{IPR and the weight norm ratio co-evolve smoothly through the
        transition} (row~d), confirming that feature formation (IPR) and
        layer balancing ($\|W_2\|/\|W_1\|$) proceed in lockstep.
\end{itemize}

Table~\ref{tab:tonset} quantifies the onset-to-grokking gap across all
conditions.  The key pattern: heterogeneity has a modest effect on
$T_{\mathrm{onset}}$ (the onset delay is 5--17\%) but a much larger effect on
the transition width $T_{\mathrm{grok}} - T_{\mathrm{onset}}$ (up to 71\%
wider).  The bottleneck is in the \emph{acceleration phase}, not the ignition.

\begin{table}[htbp]
\centering
\small
\caption{\textbf{Generalization onset and transition width} (seed~42).
$T_{\mathrm{onset}}$: first step with test accuracy $>$5\%.
Gap $= T_{\mathrm{grok}} - T_{\mathrm{onset}}$.}
\label{tab:tonset}
\begin{tabular}{llrrrr}
\toprule
$\alpha$ & $\alpha_{\mathrm{dir}}$ & $T_{\mathrm{onset}}$ &
$T_{\mathrm{grok}}$ & Gap & Gap widening \\
\midrule
0.50 & 1000 (IID) &  5{,}765 &  7{,}565 &  1{,}800 & --- \\
0.50 & 0.01       &  5{,}690 &  7{,}680 &  1{,}990 & $+$11\% \\
\addlinespace
0.30 & 1000 (IID) &  7{,}880 & 12{,}790 &  4{,}910 & --- \\
0.30 & 0.01       &  8{,}300 & 14{,}030 &  5{,}730 & $+$17\% \\
\addlinespace
0.25 & 1000 (IID) & 12{,}590 & 24{,}675 & 12{,}085 & --- \\
0.25 & 0.01       & 14{,}725 & 35{,}370 & 20{,}645 & $+$71\% \\
\bottomrule
\end{tabular}
\end{table}

\begin{figure}[htbp]
  \centering
  \includegraphics[width=\textwidth]{figures/exp3_grokking_anatomy.png}
  \caption{\textbf{Grokking anatomy} (seed~42).  Four conditions spanning
  comfortable ($\alpha{=}0.3$) to near-boundary ($\alpha{=}0.25$), under IID
  and extreme non-IID.  Green dotted line: $T_{\mathrm{onset}}$ (first
  $>$5\% test accuracy); red dashed: $T_{\mathrm{grok}}$ ($>$95\%
  permanently).  Rows: (a)~test accuracy, (b)~generalization gap,
  (c)~mean client drift, (d)~IPR and $\|W_2\|/\|W_1\|$ ratio.}
  \label{fig:anatomy}
\end{figure}


\subsubsection{Weight norm dynamics}

The Gromov mechanism requires weight norms to grow past a critical threshold
for the periodic Fourier features to emerge.
Figure~\ref{fig:weight-norms} reveals that heterogeneity \textbf{slows weight
growth in both layers}, providing a direct mechanistic explanation for the
grokking delay:

\begin{itemize}
  \item Under extreme non-IID ($\alpha_{\mathrm{dir}}{=}0.01$), layer-1 and
        layer-2 weight norms grow ${\sim}20$--$30\%$ more slowly than under IID
        during the pre-grokking phase.

  \item However, the \emph{final} weight magnitudes converge to the same
        values across all heterogeneity levels, confirming that the same
        Fourier solution is reached regardless of the data distribution.

  \item The weight norm ratio $\|W_2\|/\|W_1\|$
        (Figure~\ref{fig:weight-ratio}) converges to ${\sim}0.8$--$0.9$ for
        all conditions that grok.  Under non-IID, the ratio is consistently
        lower during training---layer~2 (the output readout) grows more
        slowly relative to layer~1 (the feature extractor), suggesting that
        aggregation noise from mismatched client label distributions
        differentially impedes the output layer.
\end{itemize}

\begin{figure}[htbp]
  \centering
  \includegraphics[width=\textwidth]{figures/exp3_weight_norm_evolution.png}
  \caption{\textbf{Weight norm evolution} (seed~42).  Top: layer-1 norm;
  bottom: layer-2 norm.  Non-IID slows growth in both layers, with a
  larger effect on layer~2.  Final norms converge regardless of heterogeneity.}
  \label{fig:weight-norms}
\end{figure}

\begin{figure}[htbp]
  \centering
  \includegraphics[width=\textwidth]{figures/exp3_weight_norm_ratio.png}
  \caption{\textbf{Weight norm ratio} ($\|W_2\|/\|W_1\|$) over training
  (seed~42).  All grokking runs converge to the same ratio $\approx 0.8$--$0.9$.
  Under non-IID, the ratio is suppressed during training, indicating
  heterogeneity preferentially slows the output layer.}
  \label{fig:weight-ratio}
\end{figure}


\subsubsection{Client drift and feature quality}

Figure~\ref{fig:drift-dynamics} shows mean client drift and client weight
divergence across heterogeneity levels.  Two distinct phenomena are visible:

\begin{itemize}
  \item \textbf{Drift magnitude scales with heterogeneity} throughout
        training (top row): $\alpha_{\mathrm{dir}}{=}0.01$ produces
        $3$--$5\times$ higher drift than IID.

  \item \textbf{Weight divergence spikes at the grokking transition}
        (bottom row).  This spike appears in all conditions and coincides
        precisely with the test accuracy jump.  It suggests that grokking
        is a cooperative phase transition that temporarily destabilizes
        client consensus---clients briefly explore different regions of the
        solution landscape before snapping into agreement on the Fourier
        solution.
\end{itemize}

Despite high drift, Figure~\ref{fig:drift-ipr} shows that \textbf{feature
quality at grokking is independent of drift level}.  The IPR at
$T_{\mathrm{grok}}$ clusters tightly around $0.13$--$0.16$ for all grokked
runs (across all $\alpha$, $\alpha_{\mathrm{dir}}$, and seeds), with no
correlation to the mean client drift experienced during training.  Drift
delays grokking but does not degrade the solution.

\begin{figure}[htbp]
  \centering
  \includegraphics[width=\textwidth]{figures/exp3_client_drift_dynamics.png}
  \caption{\textbf{Client drift dynamics} (seed~42).  Top: mean client drift;
  bottom: client weight divergence.  Drift scales with heterogeneity;
  weight divergence shows a characteristic spike at the grokking transition.}
  \label{fig:drift-dynamics}
\end{figure}

\begin{figure}[htbp]
  \centering
  \includegraphics[width=0.85\textwidth]{figures/exp3_drift_ipr_correlation.png}
  \caption{\textbf{Drift--IPR correlation at grokking.}  Each point is one
  run (all seeds).  Left: client drift vs.\ IPR at $T_{\mathrm{grok}}$;
  right: weight divergence vs.\ IPR.  Among grokked runs, IPR is tightly
  clustered regardless of drift---feature quality is independent of the
  heterogeneity experienced during training.  The low-IPR cluster
  (${\sim}0.03$--$0.06$) corresponds to $\alpha{=}0.2$ runs that did not grok.}
  \label{fig:drift-ipr}
\end{figure}


\subsubsection{Reproducibility and seed variability}

Figure~\ref{fig:per-seed} shows per-seed trajectories for four conditions
spanning comfortable and near-boundary regimes under IID and extreme non-IID:

\begin{itemize}
  \item At $\alpha{=}0.5$, all three seeds produce tightly clustered
        trajectories ($T_{\mathrm{grok}}$ varies by ${\sim}1{,}000$ steps)
        regardless of heterogeneity.

  \item At $\alpha{=}0.25$ with $\alpha_{\mathrm{dir}}{=}0.01$ (the most
        stressed condition that still groks), seed variance increases
        substantially: the fastest seed groks at ${\sim}35{,}000$ steps while
        the slowest requires ${\sim}39{,}000$.  The grokking transition
        becomes a more stochastic event near the phase boundary under high
        heterogeneity.

  \item Despite timing variability, all seeds reach 100\% test accuracy and
        converge to the same final IPR (${\sim}0.39$), confirming that the
        outcome is deterministic even when the timing is not.
\end{itemize}

\begin{figure}[htbp]
  \centering
  \includegraphics[width=\textwidth]{figures/exp3_per_seed_variability.png}
  \caption{\textbf{Per-seed variability.}  Top: test accuracy; bottom: IPR.
  Columns span from comfortable ($\alpha{=}0.5$, IID) to maximally stressed
  ($\alpha{=}0.25$, $\alpha_{\mathrm{dir}}{=}0.01$).  Timing variance
  increases near the boundary under heterogeneity, but all seeds converge to
  the same solution.}
  \label{fig:per-seed}
\end{figure}


% ---------------------------------------------------------------------------
\subsection{Summary of Findings}
\label{sec:exp3-summary}

The results from Experiment~3 support a clear narrative:

\begin{enumerate}
  \item \textbf{Grokking is robust to heterogeneity.}  Across 147 federated
        runs spanning Dirichlet, operand, and target partitions,
        \emph{every} run with $\alpha \geq 0.25$ reached 100\% test
        accuracy.  The phase boundary $\alpha_{\mathrm{crit}} \approx 0.2$
        is unchanged by heterogeneity.

  \item \textbf{Heterogeneity delays but does not prevent grokking.}
        The delay is concentrated near the phase boundary: up to 37\%
        slowdown at $\alpha{=}0.25$ under extreme non-IID
        ($\alpha_{\mathrm{dir}}{=}0.01$), but $<$1.3\% at $\alpha{=}0.5$.

  \item \textbf{The delay mechanism is weight growth retardation.}
        Non-IID data slows the growth of weight norms in both layers,
        particularly layer~2, delaying the point at which the Fourier
        features become sufficiently strong to drive generalization.

  \item \textbf{Structured partitions do not fragment Fourier learning.}
        Contrary to the initial hypothesis, operand-based partitioning
        (which fragments the input structure) produces grokking times
        within noise of IID.  The Fourier solution is learned through
        aggregation regardless of how the input space is divided across
        clients.

  \item \textbf{Client drift does not degrade the final solution.}  High
        drift during training delays grokking but has no effect on the
        quality of the learned features (IPR at convergence is identical
        across all conditions).

  \item \textbf{The grokking transition is a cooperative phase transition.}
        The characteristic spike in client weight divergence at
        $T_{\mathrm{grok}}$, the universal drift-peak-then-decay pattern,
        and the lockstep co-evolution of IPR and weight norm ratios all
        point to grokking as a sharp, coordinated transition in parameter
        space that occurs simultaneously across all clients despite their
        heterogeneous local data.
\end{enumerate}

\end{document}
